\chapter{Lo strato di ciclo: il motore diventa uno strumento per RDE}
\label{ch:cycle}

Finora l'ugello aveva un solo punto di funzionamento: una pressione
di camera, una temperatura di camera, un progetto. Un motore a
detonazione rotante non ne offre uno solo. Questo capitolo insegna
all'ottimizzatore a progettare rispetto a un intero ciclo periodico
di condizioni di camera --- e dimostra poi che lo ha fatto
correttamente, facendogli riprodurre alla cieca due risposte che la
teoria già conosce e che un'implementazione con scorciatoie non
potrebbe trovare.

\section{Ottimizzazione mediata sul ciclo e i suoi oracoli}

\subsection{Che cosa occorre sapere prima}

Un \emph{motore a detonazione rotante} (RDE) brucia il proprio
propellente in una camera anulare attorno alla quale una o più
\emph{onde di detonazione} --- fronti di combustione supersonici, un
urto con la fiamma agganciata a esso --- corrono di continuo, tipicamente
qualche migliaio di volte al secondo, con propellente fresco che
riempie lo spazio dietro ciascuna onda. La conseguenza per l'ugello è
il punto di questo capitolo: una sezione di ugello non vede mai
condizioni stazionarie. Una volta per giro la pressione salta
bruscamente al passaggio dell'onda, poi cala (\emph{si svuota}) fino
all'arrivo dell'onda successiva.

Ciò a cui l'ugello risponde è lo \emph{stato di ristagno}
$(P_0, T_0)$: la pressione e la temperatura che il gas di camera
avrebbe se portato dolcemente a riposo, la condizione di serbatoio
che fissa tutto ciò che sta a valle. In un RDE è una \emph{famiglia}
periodica $(P_0(\xi), T_0(\xi))$ indicizzata dalla \emph{fase di
ciclo} $\xi \in [0,1]$, un giro completo dell'onda; la \emph{misura
di fase} $\mathrm{d}\mu(\xi)$ dice quanta parte del ciclo occupa
ciascuna fase, e dunque il suo peso in una media. E
\emph{quasi-stazionario} significa che il gas attraversa l'ugello in
un tempo molto minore di quello che l'onda impiega a ritornare ---
cosicché in ogni istante il flusso è semplicemente il flusso
stazionario dello stato di camera di quell'istante. La variabile di
progetto è ovunque il rapporto di espansione
$\varepsilon = A_e/A_t$, area di uscita su area di gola
(Capitolo~\ref{ch:nozzle}).

\subsection{La domanda}

Un RDE consegna all'ugello una \emph{famiglia} periodica di stati di
camera, non uno solo. Il motore sa ottimizzare la spinta mediata sul
ciclo --- e dimostrare di farlo correttamente?

\subsection{Come funziona, per esteso}

Poiché l'espansione è quasi-stazionaria, la spinta mediata sul ciclo
è semplicemente la media in fase delle spinte stazionarie:
\begin{equation}
J_{\rm cycle}(\varepsilon) \;=\; \int_0^1
F\!\left(\varepsilon;\, P_0(\xi),\, T_0(\xi)\right)\mathrm{d}\mu(\xi).
\label{eq:jcycle}
\end{equation}
Lo strato implementa questa formula \emph{onestamente}: l'integrale è
una quadratura su un insieme finito di fasi, e ogni fase riceve il
proprio stato di ristagno costruito sulle stesse tabelle di gas, la
propria marcia completa attraverso l'ugello e la propria spinta. Da
nessuna parte è programmata una scorciatoia che sfrutti una
similitudine fra le fasi --- deliberatamente, perché il test più
difficile della teoria riguarda proprio ciò che
un'implementazione onesta deve allora riprodurre da sé.

Quel test è il teorema di collasso per la campana fissa (la sua
T3)~\cite{theory}. Dice: se sul ciclo varia solo la pressione e non
la temperatura (uno svuotamento \emph{a sola pressione}, $T_0$
fisso), l'intero problema di ciclo collassa su un classico progetto
stazionario alla pressione media --- quindi l'ottimizzazione
consapevole del ciclo deve comprare \emph{esattamente nulla}. La
ragione è una similitudine: a $T_0$ fisso e geometria fissa, scalare
$P_0$ scala ogni pressione nell'ugello dello stesso fattore e non
muove né pareti né numeri di Mach, quindi $J_{\rm cycle}$ è
esattamente il $J$ stazionario a $\langle P_0\rangle$. Un motore che
marcia ogni fase in modo indipendente non ha modo di saperlo --- a
meno che non sia corretto, nel qual caso il guadagno nullo
\emph{emerge}. Quell'emergere è il test.

Il secondo stadio pone il primo RDE simulato (posto dall'utente):
profili esponenziali in fase di $P_0$ \emph{e} $T_0$ (rapporto di
pressione picco-valle $PR = 10$, rapporto di temperatura
$TR = 1.306$, picco 3850~K, cinque fasi equispaziate nel logaritmo
della pressione). Con $T_0$ variabile il collasso non è più esatto, e
la domanda diventa quantitativa: quanto si avvicina all'ottimo di
ciclo trovato alla cieca il miglior progetto \emph{a un solo numero}
--- un ordinario progetto a punto singolo su uno stato di camera
rappresentativo --- e qual è lo stato giusto? La risposta della
teoria --- la pressione media alla temperatura \emph{pesata sulla
pressione} $\langle P T\rangle/\langle P\rangle$ --- viene
confrontata con l'alternativa della media logaritmica in un paragone
derivato, e un progetto fatto alla fase di picco di pressione
fornisce il riferimento deliberatamente sbagliato.

Costruire lo stato di ristagno di una fase su proprietà di gas reale
tabulate ha trappole proprie; le tre incontrate (e i rigettatori che
ora le presidiano) sono nell'Appendice~\ref{app:cycletraps}.

\subsection{Risultati di riferimento}

\textbf{L'oracolo del collasso a guadagno nullo: 5/5 PASS}
(\carrier{validation/a1\_cycle\_layer.py}, stadio \code{t3}).
L'ottimizzatore, cui non è stato detto nulla del teorema e che spende
cinque marce complete per ogni valutazione dell'obiettivo, atterra
sul progetto-alla-media entro la propria finestra: il guadagno che
trova rispetto a quel progetto è di $-475$~N su
$1.16\times10^{8}$~N, contro una banda derivata di
$9.6\times10^{3}$~N --- zero, alla precisione che il test può
risolvere. Il progetto deliberatamente sbagliato alla pressione di
picco perde $4.8\times10^{6}$~N, $500\times$ la banda, quindi un
guadagno reale si sarebbe visto. E la similitudine si manifesta dal
vivo: a geometria fissa le cinque pareti per fase coincidono a meno
di $8.9\times10^{-16}$ --- scalare la pressione non muove la parete,
misurato allo zero di macchina.

\textbf{L'RDE simulato: 6/6 PASS} (stadio \code{mock}). L'ottimo
trovato alla cieca cade entro il $-0.74\%$ in $\varepsilon$ rispetto
al miglior progetto a un solo numero --- che è appunto la pressione
media alla temperatura pesata sulla pressione
$\langle PT\rangle/\langle P\rangle$, battendo il riferimento a media
logaritmica in un paragone derivato (la lezione della media pesata
della teoria, confermata a livello di motore); la penalità di spinta
di quel progetto a un solo numero è $2.5\times10^{-6}$ relativa
(secondo ordine, come previsto), mentre il progetto di picco perde il
4\% di $J$. La lettura pratica: per questa classe di cicli uno stato
ben scelto basta, e la teoria dice quale.

\begin{figure}[htbp]
\centering
\includegraphics[width=\textwidth]{fig_cycle.pdf}
\caption{Lo strato di ciclo, misurato. Assi in entrambi i pannelli:
rapporto di espansione $\varepsilon$ [--] (area di uscita su area di
gola, la variabile di progetto) contro la spinta mediata sul ciclo
dell'Equazione~\eqref{eq:jcycle}, riportata come deficit sotto il
massimo della curva stessa, in kN (zero è il miglior progetto sulla
curva; il livello assoluto è $J_{\rm cycle} \approx 116$~MN). Ogni
punto costa cinque marce complete di ugello, una per fase; le linee
verticali segnano i progetti candidati. (a)~Svuotamento a sola
pressione: l'ottimo di ciclo trovato alla cieca
($\varepsilon = 5.1685$) cade sul classico progetto alla pressione
media ($\varepsilon = 5.1506$) entro la finestra derivata --- il
collasso a guadagno nullo, riprodotto da un motore cui non era mai
stato detto che dovesse accadere. (b)~RDE simulato, con pressione
\emph{e} temperatura variabili: l'ottimo alla cieca
($\varepsilon = 5.1395$) e il progetto a un solo numero allo stato
medio pesato sulla pressione ($\varepsilon = 5.1778$) differiscono
dello $0.74\%$ in $\varepsilon$ e di $2.5\times10^{-6}$ in $J$; il
progetto alla pressione di picco ($\varepsilon = 10.29$) perde il 4\%
e cade molto fuori scala. Si noti quanto sono piatte le curve vicino
ai loro picchi: la storia vive nelle ultime cifre, ed è per questo
che le soglie sono derivate anziché stimate a occhio.}
\label{fig:cycle}
\end{figure}

\section{Il compromesso mediato, su un misuratore}
\label{sec:votes}

\subsection{La domanda}

L'affermazione più citabile della teoria sugli ugelli per
RDE~\cite{theory} è la condizione di spigolo mediata: ``nessuna fase
soddisfa la propria condizione di spigolo; la media di ciclo sì''.
Si può misurare?

\subsection{Come funziona, per esteso}

Si richiami la condizione di spigolo della
Sezione~\ref{sec:corner}: l'affermazione esatta che la parete di un
ugello si ferma nel punto giusto, il cui termine dominante è lo
scarto fra la pressione che il gas raggiunge all'uscita e la
pressione ambiente in cui scarica. All'ottimo di ciclo ogni fase
$\xi$ esprime un \emph{voto} attraverso quello stesso misuratore: il
proprio scarto di pressione di uscita $R(\xi) = p_e(\xi) - p_a$. Una
fase con $R > 0$ lascia l'ugello ancora spingendo più forte
dell'ambiente --- \emph{sotto-espansa}, vota ``allungami''. Una fase
con $R < 0$ si è espansa sotto l'ambiente --- \emph{sovra-espansa},
vota ``accorciami''. L'affermazione è che l'ottimo non soddisfa
nessuna fase individualmente: è un compromesso mediato in cui la
media pesata con $\mu$ dei voti è nulla mentre nessun singolo voto
lo è.

Il misuratore legge i voti dalle esecuzioni di ciclo memorizzate, e
ogni soglia è derivata. La banda di equilibrio attorno allo zero
viene dalla curvatura dell'obiettivo stesso e dall'ampiezza della
finestra in cui l'ottimizzatore si è fermato --- quanto lontano
dallo zero la media potrebbe legittimamente trovarsi. Ogni voto di
fase deve poi individualmente superare quella banda; altrimenti
``nessuna fase soddisfa la propria condizione'' sarebbe
un'asserzione anziché una misura.

\subsection{Risultati di riferimento}

\textbf{8/8 PASS su entrambi i cicli}
(\carrier{validation/a1\_cycle\_corner.py}, un'analisi delle
esecuzioni di ciclo memorizzate --- 0.4~s). All'ottimo dell'RDE
simulato i voti delle cinque fasi vanno da $+7.9\times10^{5}$ fino a
$-5.1\times10^{5}$~Pa --- monotoni in fase, con esattamente un
cambio di segno, ciascuno individualmente 2.8--3.2$\times$ la banda
di equilibrio e 14--90$\times$ la media raggiunta --- mentre la loro
\emph{media} sta a zero entro la banda derivata. Sul progetto
sbagliato (alla pressione di picco) il misuratore si inchioda a
$-4.6\times10^{5}$~Pa: tutte le fasi sovra-espanse, nessun
compromesso. Un progetto sbagliato è visibile a colpo d'occhio.

\begin{figure}[htbp]
\centering
\includegraphics[width=0.85\textwidth]{fig_votes.pdf}
\caption{Il teorema su un quadrante. Asse orizzontale: fase di ciclo
$\xi$ [--], un giro dell'onda dal passaggio alla coda dello
svuotamento. Asse verticale: il voto di spigolo di ciascuna fase
$p_e - p_a$ [bar], la quantità di cui la sua pressione di uscita
supera (positivo: sotto-espansa, ``allungami'') o non raggiunge
(negativo: sovra-espansa, ``accorciami'') la pressione ambiente. Blu:
i voti all'ottimo di ciclo --- grandi e di entrambi i segni, eppure
la loro media (blu tratteggiato) sta dentro la banda ombreggiata di
equilibrio allo zero. È la condizione di spigolo mediata, misurata:
nessuna fase è individualmente soddisfatta, il ciclo sì. Grigio: gli
stessi voti sul progetto alla pressione di picco --- tutte le fasi
sovra-espanse, la media (grigio tratteggiato) molto sotto lo zero: un
progetto sbagliato, a colpo d'occhio.}
\label{fig:votes}
\end{figure}
